% ========== 1.5 PROJECT SCOPE ==========

\clearpage
\section{1.5 Project Scope}
\addcontentsline{toc}{section}{1.5 Project Scope}
\label{sec:scope}

\lettrine{T}{his project} encompasses the design, architecture, and conceptual framework of Symphony as an AI-First Development Environment. The scope defines clear boundaries between what the project includes and excludes, ensuring focused development efforts and realistic deliverables within the graduation project timeline.

\subsection*{1.5.1 In Scope}
\addcontentsline{toc}{subsection}{1.5.1 In Scope}

\textbf{Core Architecture and Design}

The project includes complete architectural design of Symphony's microkernel system, the Dual Ensemble Architecture (The Pit and The Grand Stage), and the intelligent orchestration framework. This encompasses system component specifications, interaction protocols, and integration patterns that enable AI-first development workflows.

\textbf{Orchestration System}

The Conductor's reinforcement learning architecture, including the Function Quest Game training methodology, reward system design, and multi-agent coordination mechanisms. This includes the Melody workflow system for visual composition and execution of AI-driven development tasks.

\textbf{Extension System Framework}

The Orchestra Kit architecture for extension development, distribution, and lifecycle management. This includes the Player system for extension governance, security models for sandboxed execution, and the marketplace infrastructure for community-driven innovation.

\textbf{Infrastructure Components}

Design specifications for The Pit's five core extensions: Pool Manager for AI model lifecycle management, DAG Tracker for workflow dependency resolution, Artifact Store for institutional memory, Arbitration Engine for resource allocation, and Stale Manager for system optimization.

\textbf{User Interface Design}

Complete UI/UX design for the IDE interface, including the Harmony Board for visual workflow composition, extension marketplace integration, and the three operational modes (Maestro, Virtuoso, Solo). Website design for project presentation and documentation.

\textbf{System Integration}

Integration patterns between Symphony's components, IPC communication protocols, data flow architectures, and the H²A² (Harmonic Hexagonal Actor Architecture) implementation strategy.

\subsection*{1.5.2 Out of Scope}
\addcontentsline{toc}{subsection}{1.5.2 Out of Scope}

\textbf{Full Production Implementation}

Complete production-ready implementation of all system components is beyond the graduation project scope. The focus is on architectural design, proof-of-concept implementations, and demonstrating feasibility rather than market-ready software.

\textbf{Large Language Model Training}

Training custom large language models (LLMs) for code generation, completion, or understanding is excluded. The project utilizes existing LLM APIs (GPT, Claude, etc.) as external services. However, the project includes training the Conductor's reinforcement learning model from scratch using the Function Quest Game methodology for orchestration intelligence.


\textbf{Enterprise Features}

Advanced enterprise features such as multi-tenant architecture, enterprise-grade authentication systems, compliance frameworks, and large-scale deployment infrastructure are not included in this graduation project scope.

\textbf{Cross-Platform Mobile Support}

Mobile application development for iOS and Android platforms is excluded. The project focuses on desktop and web-based development environments.

\textbf{Extensive Extension Library}

Building a comprehensive library of community extensions is beyond the project timeline. The focus is on the extension framework and demonstrating extensibility through representative examples.

\textbf{Long-Term Performance Optimization}

While performance considerations are included in the design, extensive performance tuning, benchmarking across diverse hardware configurations, and production-scale load testing are outside the graduation project scope.

% \subsection*{1.5.3 Project Boundaries}
% \addcontentsline{toc}{subsection}{1.5.3 Project Boundaries}

% \textbf{Technical Boundaries}

% The project demonstrates Symphony's architecture through targeted implementations that prove key concepts. Focus areas include the orchestration system's core functionality, extension framework viability, and user interface design. Implementation depth varies by component based on research value and demonstration requirements.

% \textbf{Timeline Boundaries}

% Development spans the academic graduation project timeline, with milestones aligned to academic requirements. The project prioritizes architectural innovation and design validation over feature completeness.

% \textbf{Resource Boundaries}

% The project operates within academic resource constraints, utilizing available computing resources, open-source tools, and existing AI model APIs. Commercial licensing, cloud infrastructure costs, and specialized hardware are limited to what is accessible within academic budgets.
